% =====================================================================
% COMECE AQUI.
%
% Regra de ouro: mude apenas o texto que está entre as chaves { }.
% Não apague o comando (a palavra que começa com \), nem as chaves.
%
%   \autor{Nome do aluno}
%   ^^^^^^ comando        ^ conteúdo: é só isto que você troca
%
% Estes dados são reaproveitados na capa, na folha de rosto, no resumo,
% nas notas de rodapé e nas propriedades do PDF. Preencher aqui basta;
% você não precisa repetir seu nome em nenhum outro arquivo.
% =====================================================================


% ---------------------------------------------------------------------
% 1. CURSO E CAMPUS  (obrigatório)
%
% Definidos uma única vez e reaproveitados em todos os lugares abaixo.
% Se você trocar o curso aqui, a capa, a folha de rosto e a nota de
% rodapé mudam juntas.
% ---------------------------------------------------------------------
\newcommand{\meucurso}{Análise e Desenvolvimento de Sistemas}
\newcommand{\meucampus}{Campus Pedro II}


% ---------------------------------------------------------------------
% 2. IDENTIFICAÇÃO DO TRABALHO  (obrigatório)
% ---------------------------------------------------------------------

% Título do artigo. Se houver subtítulo, separe-o por dois-pontos.
\titulo{Título do trabalho: subtítulo (se houver)}

% Seu nome completo, como deve aparecer na capa.
\autor{Nome do aluno}

% Cidade e estado que aparecem no rodapé da capa e da folha de rosto.
\local{Pedro II, Piauí}

% Ano de conclusão. Aparece logo abaixo da cidade.
\data{2026}

% Nome do curso para o abnTeX2 (usado na capa de CD, se você ativá-la).
\nomedocurso{\meucurso}


% ---------------------------------------------------------------------
% 3. INSTITUIÇÃO  (obrigatório)
%
% ATENÇÃO: aqui o \\ no fim da linha é obrigatório — ele força a quebra
% de linha no cabeçalho da capa. Mantenha um \\ ao fim de cada linha,
% menos na última. (No texto corrido dos capítulos, ao contrário, você
% NÃO deve usar \\ para quebrar linhas.)
%
% A terceira linha traz o nome do curso como ele aparece no diploma
% ("TECNOLOGIA EM ..."). Cursos de bacharelado ou licenciatura devem
% trocar esse prefixo, ou apagá-lo.
% ---------------------------------------------------------------------
\instituicao{%
  INSTITUTO FEDERAL DE EDUCAÇÃO, CIÊNCIA E TECNOLOGIA DO PIAUÍ\\
  \MakeUppercase{\meucampus}\\
  TECNOLOGIA EM \MakeUppercase{\meucurso}}


% ---------------------------------------------------------------------
% 4. TEXTO DA FOLHA DE ROSTO  (obrigatório)
%
% \tipotrabalho vai só nas propriedades do PDF (não aparece na página).
% \preambulo é o parágrafo que fica no lado direito da folha de rosto,
% explicando a que o trabalho se destina. Cuidado: este \preambulo não
% tem relação com o "preâmbulo" do LaTeX (o arquivo estrutura/preambulo
% .tex, que carrega os pacotes). São coisas diferentes com o mesmo nome.
% ---------------------------------------------------------------------
\tipotrabalho{Trabalho de Conclusão de Curso (Artigo Científico)}

\preambulo{Trabalho de Conclusão de Curso (artigo científico) apresentado
como exigência parcial para obtenção do diploma do Curso de \meucurso\ do
Instituto Federal de Educação, Ciência e Tecnologia do Piauí, \meucampus.}


% ---------------------------------------------------------------------
% 5. ORIENTAÇÃO  (orientador obrigatório; coorientador opcional)
%
% Escreva a titulação junto com o nome, por exemplo "Prof. Dr. Fulano".
% O rótulo "Orientador:" / "Coorientador:" é inserido automaticamente.
% ---------------------------------------------------------------------
\orientador{Prof. Dr. Nome do orientador}

% Sem coorientação? Deixe as chaves vazias: a linha some da folha de rosto.
\coorientador{}


% ---------------------------------------------------------------------
% 6. BANCA EXAMINADORA  (opcional)
%
% Estes nomes só aparecem se você ativar uma folha de aprovação em
% estrutura/pre-textual.tex. Enquanto ela estiver desativada, preencher
% aqui não muda nada no PDF.
%
%   \imprimirfolhadeaprovacao             -> orientador + membro um + dois
%   \imprimirfolhadeaprovacaoduascolunas  -> os quatro nomes, em 2 colunas
%
% Tem só dois avaliadores? Use a primeira versão e deixe \membrotres vazio.
% ---------------------------------------------------------------------
\membroum{Prof. Me. Nome do primeiro membro da banca}
\membrodois{Prof. Me. Nome do segundo membro da banca}
\membrotres{Prof. Dr. Nome do terceiro membro da banca}

% Instituição impressa sob cada assinatura da banca. Descomente e edite
% apenas se a banca for de outra instituição.
% \renewcommand{\instituicaodabanca}{Instituto Federal do Piauí (IFPI)}


% ---------------------------------------------------------------------
% 7. DATA DE APRESENTAÇÃO  (obrigatório)
%
% CUIDADO: esta data já aparece no PDF, no rodapé da página do abstract,
% mesmo com a folha de aprovação desativada. Se você não trocar, sai
% "Data de aprovação: Dia de mês de ano" no trabalho entregue.
% Exemplo de preenchimento: \dataapresentacao{20 de novembro de 2026}
% ---------------------------------------------------------------------
\dataapresentacao{Dia de mês de ano}


% ---------------------------------------------------------------------
% 8. NOTAS DE RODAPÉ DOS AUTORES  (obrigatório)
%
% Aparecem como nota de rodapé na primeira página, junto dos nomes.
% Repare que aqui há DOIS pares de chaves: o primeiro é o nome do
% comando, o segundo é o texto. Troque só o texto do segundo par.
%
%   \newcommand{\notadoautor}{este texto aqui}
%                ^nome^      ^--- o que você edita ---^
%
% Troque também os endereços: não deixe example.com no trabalho final.
% ---------------------------------------------------------------------
\newcommand{\notadoautor}{Graduando(a) em \meucurso\ pelo IFPI.
Contato: \url{aluno@example.com}.}

\newcommand{\notadoorientador}{Professor(a) do Instituto Federal do Piauí.
Contato: \url{orientador@example.com}.}


% ---------------------------------------------------------------------
% 9. PALAVRAS-CHAVE  (obrigatório)
%
% Texto simples, sem comandos de formatação, separado por ponto e vírgula
% e terminado em ponto final. \keywords é a versão em inglês, usada no
% abstract; traduza as mesmas palavras, não escolha outras.
% ---------------------------------------------------------------------
\newcommand{\palavraschave}{modelo; artigo; pesquisa.}
\newcommand{\keywords}{template; article; research.}
